\section{Market Microstructure}\label{sec:microstructure}

Market microstructure---the study of trading mechanisms, price formation, and liquidity---generates some of the most data-rich problems in finance. Limit order books (LOBs) produce millions of events per day per instrument, creating both opportunities and challenges for ML methods.

\subsection{Limit Order Book Modeling}

\subsubsection{Mid-Price Prediction}

Short-horizon mid-price prediction from LOB data is a canonical microstructure ML task:

\begin{itemize}[leftmargin=*]
    \item \citet{Sirignano2019} train deep learning models on LOB data from hundreds of stocks simultaneously, finding that a universal model (trained across stocks) outperforms stock-specific models, suggesting common microstructure patterns across assets.

    \item \citet{Zhang2019LOB} apply LSTMs to LOB snapshots and demonstrate that temporal patterns in the order flow carry predictive information beyond the current book state.

    \item \citet{Kolm2023} benchmark CNNs, LSTMs, transformers, and tree-based methods on a standardized LOB dataset (FI-2010, \citealt{Ntakaris2018}), finding that transformers with positional encoding achieve the best classification accuracy for directional prediction at horizons of 10--100 events.
\end{itemize}

\subsubsection{Full Distribution Modeling}

Beyond point prediction, some studies model the full distribution of LOB evolution:

\begin{itemize}[leftmargin=*]
    \item \citet{Cont2021} use neural point processes to model the arrival of limit orders, market orders, and cancellations as a marked point process, capturing the clustering and self-exciting nature of order flow.

    \item \citet{Shi2022} apply normalizing flows to model the joint distribution of multi-level LOB states, enabling realistic scenario generation for execution simulation.
\end{itemize}

\subsection{Optimal Execution}

Optimal execution---minimizing the market impact of large orders---is a natural RL application:

\begin{itemize}[leftmargin=*]
    \item \citet{Ning2021} compare RL-based execution strategies (DQN, PPO) against the Almgren-Chriss analytical benchmark \citep{Almgren2001} and find that RL agents learn to exploit intraday liquidity patterns, achieving lower implementation shortfall.

    \item \citet{Fang2021} use multi-agent RL to model the interaction between a large executor and adaptive market makers, capturing the game-theoretic aspects of execution that single-agent models miss.

    \item \citet{Schnaubelt2022} train a meta-learner that adapts execution strategies to changing market conditions (volatility regimes, liquidity shocks) without retraining.
\end{itemize}

\subsection{Market Making}

Automated market making---providing liquidity by continuously quoting bid and ask prices---has been extensively studied with RL:

\begin{itemize}[leftmargin=*]
    \item \citet{Gueant2017} provide the theoretical foundation: Avellaneda-Stoikov style models with analytical solutions that serve as benchmarks for ML approaches.

    \item \citet{Spooner2018} apply DRL (DQN and variants) to market making in a realistic LOB simulator, showing that the learned policy adapts spread and inventory management to market conditions.

    \item \citet{Gasperov2021} propose a multi-agent market making framework where multiple RL agents interact in a simulated market, producing emergent dynamics (bid-ask bounce, adverse selection) consistent with empirical microstructure.
\end{itemize}

\subsection{Data and Access Challenges}

A distinctive feature of the microstructure literature is the data access barrier. Tick-level LOB data is expensive and often proprietary:

\begin{itemize}[leftmargin=*]
    \item The FI-2010 dataset \citep{Ntakaris2018} (Helsinki Stock Exchange, 10 stocks, 10 business days) has become a de facto standard, but its small size and non-US market limit generalizability.
    \item LOBSTER (NASDAQ data) is widely used in academic studies but requires institutional subscription.
    \item Crypto exchanges (Binance, Coinbase) provide free LOB data, making cryptocurrency microstructure an increasingly popular testbed.
\end{itemize}

This data bottleneck explains the odd situation where microstructure---a domain that generates more data per day than any other in finance---is under-represented in the ML literature. Researchers who want to work on LOB modeling face a choice between the small, dated FI-2010 dataset and expensive commercial data. The growing availability of free crypto exchange data may eventually break this logjam, though cryptocurrency market microstructure has its own peculiarities.
